% GENERATED FILE. Edit canonical lessons, then run npm run generate:books.
% canonical-source-hash: fnv1a64:b5ea63eb8cbbee95
% canonical-lessons: IT-C371-rilevante, IT-C371-cruciale, IT-C371-essenziale, IT-C371-fondamentale, IT-C371-premessa

\chapter{Ciò che conta}
\label{ch:it-cio-che-conta}
\hlchaptermodality{371}

\begin{chapteropening}
\textbf{By the end of this chapter:} \emph{I can say what matters in an argument: what is relevant, crucial, essential or fundamental, and the premise I start from.}

Complete the last lesson of chapter 371: I can say what matters in an argument: what is relevant, crucial, essential or fundamental, and the premise I start from.
\end{chapteropening}

\section[rilevante]{rilevante — relevant, significant}

\label{lesson:IT-C371-rilevante}



\begin{quote}
Before the new one: say the Italian for the reform, then the Italian for the member of parliament.
\end{quote}

\subsection*{You'll want to know: rilevante}
\textbf{rilevante} — \textquotedblleft{}relevant, significant\textquotedblright{}.

\begin{quote}
Questo dettaglio non è rilevante.
\end{quote}

\begin{quote}
\emph{This detail is not relevant.}
\end{quote}

It also means considerable: \emph{una somma rilevante}.

\begin{cousinweb}
\emph{rilevante} is the present participle of \emph{rilevare} \textquotedblleft{}to raise, to point out\textquotedblright{}, from Latin \emph{relevāre} \textquotedblleft{}to lift up again\textquotedblright{} (Lewis \& Short), built on \emph{levāre} \textquotedblleft{}to raise\textquotedblright{}: what stands out.
\end{cousinweb}

\begin{grammarlens}[title={What you've built}]
One more word for telling, explaining and describing at B1.
\end{grammarlens}

\subsection*{Your turn}
\begin{itemize}
\raggedright
  \item \emph{Say it:} \emph{rilevante}
  \item \emph{Say it:} \emph{rilevante}, once more
  \item \emph{Say it:} a sentence of your own with \emph{rilevante}
  \item \emph{Recall:} say the Italian for the reform, then the Italian for the member of parliament, then say \emph{rilevante} again
\end{itemize}

\begin{tcolorbox}[breakable,colback=teal!4,colframe=teal!35!black,title={Before you move on}]
What does rilevante mean? (\textquotedblleft{}Relevant, significant\textquotedblright{}.) Say it once more.
\end{tcolorbox}
\section[cruciale]{cruciale — crucial}

\label{lesson:IT-C371-cruciale}



\begin{quote}
Before the new one: say the Italian for the member of parliament, then the Italian for relevant.
\end{quote}

\subsection*{You'll want to know: cruciale}
\textbf{cruciale} — \textquotedblleft{}crucial\textquotedblright{}.

\begin{quote}
Il primo anno di lavoro è cruciale.
\end{quote}

\begin{quote}
\emph{The first year of work is crucial.}
\end{quote}

\emph{Un momento cruciale} — a turning point.

\begin{cousinweb}
\emph{cruciale} is usually traced to French \emph{crucial}, built on Latin \emph{crux, crucis} \textquotedblleft{}a cross\textquotedblright{} (Lewis \& Short): the point where two ways cross and you must choose one.
\end{cousinweb}

\begin{grammarlens}[title={What you've built}]
One more word for telling, explaining and describing at B1.
\end{grammarlens}

\subsection*{Your turn}
\begin{itemize}
\raggedright
  \item \emph{Say it:} \emph{cruciale}
  \item \emph{Say it:} \emph{cruciale}, once more
  \item \emph{Say it:} a sentence of your own with \emph{cruciale}
  \item \emph{Recall:} say the Italian for the member of parliament, then the Italian for relevant, then say \emph{cruciale} again
\end{itemize}

\begin{tcolorbox}[breakable,colback=teal!4,colframe=teal!35!black,title={Before you move on}]
What does cruciale mean? (\textquotedblleft{}Crucial\textquotedblright{}.) Say it once more.
\end{tcolorbox}
\section[essenziale]{essenziale — essential}

\label{lesson:IT-C371-essenziale}



\begin{quote}
Before the new one: say the Italian for relevant, then the Italian for crucial.
\end{quote}

\subsection*{You'll want to know: essenziale}
\textbf{essenziale} — \textquotedblleft{}essential\textquotedblright{}.

\begin{quote}
L'acqua è essenziale per la vita.
\end{quote}

\begin{quote}
\emph{Water is essential for life.}
\end{quote}

\emph{L'essenziale} — the main thing. \emph{Uno stile essenziale} — a plain style.

\begin{cousinweb}
\emph{essenziale} is a learned borrowing from Late Latin \emph{essentiālis} \textquotedblleft{}of the essence\textquotedblright{} (EtymDB), from \emph{essentia} \textquotedblleft{}being, essence\textquotedblright{}, built on \emph{esse} \textquotedblleft{}to be\textquotedblright{}.
\end{cousinweb}

\begin{grammarlens}[title={What you've built}]
One more word for telling, explaining and describing at B1.
\end{grammarlens}

\subsection*{Your turn}
\begin{itemize}
\raggedright
  \item \emph{Say it:} \emph{essenziale}
  \item \emph{Say it:} \emph{essenziale}, once more
  \item \emph{Say it:} a sentence of your own with \emph{essenziale}
  \item \emph{Recall:} say the Italian for relevant, then the Italian for crucial, then say \emph{essenziale} again
\end{itemize}

\begin{tcolorbox}[breakable,colback=teal!4,colframe=teal!35!black,title={Before you move on}]
What does essenziale mean? (\textquotedblleft{}Essential\textquotedblright{}.) Say it once more.
\end{tcolorbox}
\section[fondamentale]{fondamentale — fundamental, vital}

\label{lesson:IT-C371-fondamentale}



\begin{quote}
Before the new one: say the Italian for crucial, then the Italian for essential.
\end{quote}

\subsection*{You'll want to know: fondamentale}
\textbf{fondamentale} — \textquotedblleft{}fundamental, vital\textquotedblright{}.

\begin{quote}
Per questo lavoro l'inglese è fondamentale.
\end{quote}

\begin{quote}
\emph{For this job English is vital.}
\end{quote}

Very frequent in speech: \emph{è fondamentale che} + subjunctive.

\begin{cousinweb}
\emph{fondamentale} is built on \emph{fondamento} \textquotedblleft{}foundation\textquotedblright{}, inherited from Latin \emph{fundāmentum} \textquotedblleft{}a foundation\textquotedblright{} (Lewis \& Short), from \emph{fundus} \textquotedblleft{}bottom, ground\textquotedblright{}.
\end{cousinweb}

\begin{grammarlens}[title={What you've built}]
One more word for telling, explaining and describing at B1.
\end{grammarlens}

\subsection*{Your turn}
\begin{itemize}
\raggedright
  \item \emph{Say it:} \emph{fondamentale}
  \item \emph{Say it:} \emph{fondamentale}, once more
  \item \emph{Say it:} a sentence of your own with \emph{fondamentale}
  \item \emph{Recall:} say the Italian for crucial, then the Italian for essential, then say \emph{fondamentale} again
\end{itemize}

\begin{tcolorbox}[breakable,colback=teal!4,colframe=teal!35!black,title={Before you move on}]
What does fondamentale mean? (\textquotedblleft{}Fundamental, vital\textquotedblright{}.) Say it once more.
\end{tcolorbox}
\section[la premessa]{la premessa — the premise; the preface, the introduction}

\label{lesson:IT-C371-premessa}



\begin{quote}
Before the new one: say the Italian for essential, then the Italian for fundamental.
\end{quote}

\subsection*{You'll want to know: la premessa}
\textbf{la premessa} — \textquotedblleft{}the premise; the preface, the introduction\textquotedblright{}.

\begin{quote}
Faccio una premessa: non sono un esperto.
\end{quote}

\begin{quote}
\emph{Let me say first: I am not an expert.}
\end{quote}

\emph{Fare una premessa} — to say something first. In a book it is the preface.

\begin{cousinweb}
\emph{premessa} is the past participle of \emph{premettere} \textquotedblleft{}to put first\textquotedblright{}, from Latin \emph{praemittere} \textquotedblleft{}to send ahead\textquotedblright{} (Lewis \& Short), built from \emph{prae} \textquotedblleft{}before\textquotedblright{} and \emph{mittere} \textquotedblleft{}to send\textquotedblright{}.
\end{cousinweb}

\begin{grammarlens}[title={What you've built}]
One more word for telling, explaining and describing at B1.
\end{grammarlens}

\subsection*{Your turn}
\begin{itemize}
\raggedright
  \item \emph{Say it:} \emph{la premessa}
  \item \emph{Say it:} \emph{la premessa}, once more
  \item \emph{Say it:} a sentence of your own with \emph{la premessa}
  \item \emph{Recall:} say the Italian for essential, then the Italian for fundamental, then say \emph{la premessa} again
\end{itemize}

\begin{tcolorbox}[breakable,colback=teal!4,colframe=teal!35!black,title={Before you move on}]
What does la premessa mean? (\textquotedblleft{}The premise; the preface, the introduction\textquotedblright{}.) Say it once more.
\end{tcolorbox}
